\documentclass[11pt]{article}

\usepackage{amsmath,amssymb,amsthm}
\usepackage[margin=1.15in]{geometry}
\usepackage{enumitem}
\usepackage{hyperref}

\newtheorem{theorem}{Theorem}[section]
\newtheorem{lemma}[theorem]{Lemma}
\newtheorem{corollary}[theorem]{Corollary}
\newtheorem{proposition}[theorem]{Proposition}
\theoremstyle{remark}
\newtheorem{remark}[theorem]{Remark}

\newcommand{\cb}{\mathrm{cb}}
\newcommand{\id}{\mathrm{id}}
\newcommand{\CB}{\mathrm{CB}}
\newcommand{\CN}{\mathcal N}
\newcommand{\CI}{\mathcal I}
\newcommand{\tr}{\mathrm{tr}}
\newcommand{\ohat}{\widehat{\otimes}}
\newcommand{\ocheck}{\check{\otimes}}
\newcommand{\Minj}{M_{n,k}}
\newcommand{\Tnk}{T_{n,k}}

\title{A corrigendum on completely integral nuclearity and a structure theorem}
\author{Yafei Zhao}
\date{}

\begin{document}

\maketitle

\begin{abstract}
We correct Lemma 3.13 and Theorem 3.14 of Dong--Tao--Zhao,
\emph{Nuclearity and finite-representability in the system of completely
integral mapping spaces}.  The completely integral norm of the identity
on the rectangular matrix space \(M_{n,k}\) is \(nk\), not \(\sqrt{nk}\).
The latter number is obtained by using a Banach trace class norm after
identifying the representing tensor with a rank-one matrix in
\(M_{n^2,k^2}\), but the completely nuclear norm is the operator space
projective tensor norm on \(T_{n,k}\ohat M_{n,k}\).  We give the corrected
statement and prove the resulting structure theorem: an operator space is
\(\lambda\)-nuclear in the system of completely integral mapping spaces
if and only if it is a finite \(\ell_\infty\)-sum of rectangular matrix
spaces \(M_{n_i,k_i}\) with \(\sum_i n_i k_i\leq \lambda\).
\end{abstract}

\noindent\textbf{Keywords.} Operator space; completely integral mapping;
completely nuclear mapping; \(\lambda\)-nuclearity; injective operator space.

\medskip
\noindent\textbf{2020 Mathematics Subject Classification.} 46L07, 47L25.

\section{The correction}

Let \(M_{n,k}\) denote the \(n\times k\) rectangular matrix operator space
and let \(T_{n,k}=M_{n,k}^{*}\) be its trace class dual.  For operator
spaces \(V,W\), we write \(\CN(V,W)\) and \(\CI(V,W)\) for the completely
nuclear and completely integral mapping spaces, with norms \(\nu\) and
\(\iota\), respectively.

The proof of \cite[Lemma 3.13]{DTZ} identifies the tensor representing
\(\id_{M_{n,k}}\) with
\[
        \sum_{i=1}^{n}\sum_{j=1}^{k} E_{ij}^{*}\otimes E_{ij}
        \in T_{n,k}\ohat M_{n,k}.
\]
It then computes its norm as \(\sqrt{nk}\).  That computation is not a
computation of the operator space projective tensor norm.  It computes the
Banach trace class norm of the rank-one matrix obtained after viewing the
same formal sum inside \(M_{n^2,k^2}\).  These norms need not agree, and
the completely nuclear norm is the quotient/operator space projective norm
coming from \(T_{n,k}\ohat M_{n,k}\).

The corrected value is the following.

\begin{lemma}\label{lem:id-rectangular}
For every \(n,k\in \mathbb N\),
\[
        \iota(\id_{M_{n,k}})=\nu(\id_{M_{n,k}})=nk.
\]
\end{lemma}

\begin{proof}
Since \(M_{n,k}\) is finite-dimensional, the completely integral and
completely nuclear norms agree on \(\CB(M_{n,k},M_{n,k})\).  Moreover
\[
        \CN(M_{n,k},M_{n,k}) = T_{n,k}\ohat M_{n,k}
\]
isometrically.

Let \(\{E_{ij}\}\) be the canonical matrix units of \(M_{n,k}\), and let
\(\{E_{ij}^{*}\}\subset T_{n,k}\) be the biorthogonal functionals.  Then
\(\|E_{ij}\|=\|E_{ij}^{*}\|=1\), and
\[
        \id_{M_{n,k}}
        =
        \sum_{i=1}^{n}\sum_{j=1}^{k} E_{ij}^{*}\otimes E_{ij}
        \quad\hbox{in }T_{n,k}\ohat M_{n,k}.
\]
The projective tensor norm gives
\[
        \nu(\id_{M_{n,k}})
        \leq
        \sum_{i=1}^{n}\sum_{j=1}^{k}
        \|E_{ij}^{*}\|\,\|E_{ij}\|
        = nk.
\]
For the reverse inequality, use the contractive evaluation functional
\[
        \mathrm{ev}:T_{n,k}\ohat M_{n,k}\longrightarrow \mathbb C,
        \qquad
        \mathrm{ev}(f\otimes x)=f(x).
\]
Applying it to the above tensor yields
\[
        \nu(\id_{M_{n,k}})
        \geq
        |\mathrm{ev}(\id_{M_{n,k}})|
        =
        \sum_{i=1}^{n}\sum_{j=1}^{k} E_{ij}^{*}(E_{ij})
        = nk.
\]
Thus \(\nu(\id_{M_{n,k}})=nk\), and hence
\(\iota(\id_{M_{n,k}})=nk\).
\end{proof}

\begin{corollary}\label{cor:corrected-rectangular}
Let \(n,k\in\mathbb N\) and \(\lambda>0\).  Then \(M_{n,k}\) is
\(\lambda\)-nuclear in the system \((\CI(\cdot,\cdot),\iota)\) if and
only if
\[
        nk\leq \lambda.
\]
\end{corollary}

\begin{proof}
The proof of \cite[Lemma 3.12]{DTZ} gives that an operator space \(V\) is
\(\lambda\)-nuclear in the system \((\CI(\cdot,\cdot),\iota)\) if and
only if \(V\) is nuclear in the usual operator space sense and
\(\iota(\id_V)\leq \lambda\).  Since \(M_{n,k}\) is injective and hence
nuclear in the usual sense, the assertion follows from
Lemma~\ref{lem:id-rectangular}.
\end{proof}

\begin{remark}
Corollary~\ref{cor:corrected-rectangular} replaces
\cite[Theorem 3.14]{DTZ}.  In particular, the threshold for
\(\lambda\)-nuclearity of \(M_{n,k}\) in the completely integral system is
the dimension \(nk\), not \(\sqrt{nk}\).
\end{remark}

\section{Finite direct sums}

We next record the corresponding computation for finite injective
operator spaces.  It is the point at which the correction changes the
quantitative form of the structure theorem.

\begin{lemma}\label{lem:id-sum}
Let \(n_i,k_i\in\mathbb N\) for \(1\leq i\leq p\), and set
\[
        L=M_{n_1,k_1}\oplus_\infty\cdots\oplus_\infty M_{n_p,k_p}.
\]
Then
\[
        \iota(\id_L)=\nu(\id_L)=\sum_{i=1}^{p} n_i k_i=\dim L.
\]
\end{lemma}

\begin{proof}
The space \(L\) is finite-dimensional, so \(\iota(\id_L)=\nu(\id_L)\).
Let \(J_i:M_{n_i,k_i}\to L\) and
\(P_i:L\to M_{n_i,k_i}\) be the canonical complete isometric inclusions
and complete contractive projections.  Since
\[
        \id_L=\sum_{i=1}^{p} J_i\circ \id_{M_{n_i,k_i}}\circ P_i,
\]
the operator ideal property of \(\nu\), together with
Lemma~\ref{lem:id-rectangular}, gives
\[
        \nu(\id_L)\leq \sum_{i=1}^{p} n_i k_i.
\]
The evaluation functional on \(L^{*}\ohat L\) is contractive.  Evaluating
the tensor corresponding to \(\id_L\) gives \(\dim L\).  Hence
\[
        \nu(\id_L)\geq \dim L=\sum_{i=1}^{p} n_i k_i,
\]
and the equality follows.
\end{proof}

\section{The structure theorem}

We use the following convention.  An operator space \(V\) is
\(\lambda\)-nuclear in the system \((\CI(\cdot,\cdot),\iota)\) if there
are matrix algebras \(M_{m_\alpha}\) and linear maps
\[
        V \xrightarrow{\ \varphi_\alpha\ } M_{m_\alpha}
        \xrightarrow{\ \psi_\alpha\ } V
\]
such that \(\iota(\varphi_\alpha)\leq \lambda\),
\(\iota(\psi_\alpha)\leq \lambda\), and
\(\psi_\alpha\varphi_\alpha\to \id_V\) in the point-norm topology.  This
is the convention used in \cite{DTZ}.

\begin{theorem}\label{thm:integral-structure}
Let \(V\) be an operator space and let \(\lambda\geq 1\).  Then \(V\) is
\(\lambda\)-nuclear in the system \((\CI(\cdot,\cdot),\iota)\) if and
only if there exist \(n_i,k_i\in\mathbb N\), \(1\leq i\leq p\), such that
\[
        V
        =
        M_{n_1,k_1}\oplus_\infty\cdots\oplus_\infty M_{n_p,k_p}
        \quad\hbox{completely isometrically}
\]
and
\[
        \sum_{i=1}^{p} n_i k_i\leq \lambda.
\]
\end{theorem}

\begin{proof}
Assume first that \(V\) is \(\lambda\)-nuclear in
\((\CI(\cdot,\cdot),\iota)\).  By the characterization used in
Corollary~\ref{cor:corrected-rectangular}, \(V\) is nuclear in the usual
operator space sense and \(\iota(\id_V)\leq \lambda\).  The latter
condition forces \(V\) to be finite-dimensional.  Thus \(V^{**}=V\), and
usual nuclearity implies that \(V\) is injective.

The finite-dimensional injective operator spaces are exactly the finite
\(\ell_\infty\)-sums of rectangular matrix spaces.  Hence
\[
        V
        =
        M_{n_1,k_1}\oplus_\infty\cdots\oplus_\infty M_{n_p,k_p}
\]
completely isometrically for suitable \(n_i,k_i\).  Lemma~\ref{lem:id-sum}
then gives
\[
        \sum_{i=1}^{p} n_i k_i
        =
        \iota(\id_V)
        \leq \lambda.
\]

Conversely, suppose that \(V\) has the displayed form and that
\(\sum_i n_i k_i\leq \lambda\).  Such a finite \(\ell_\infty\)-sum is
injective, hence nuclear in the usual operator space sense.  By
Lemma~\ref{lem:id-sum},
\[
        \iota(\id_V)=\sum_{i=1}^{p}n_i k_i\leq\lambda.
\]
The same characterization therefore shows that \(V\) is
\(\lambda\)-nuclear in \((\CI(\cdot,\cdot),\iota)\).
\end{proof}

\begin{corollary}\label{cor:one-nuclear}
An operator space is nuclear in the system
\((\CI(\cdot,\cdot),\iota)\) if and only if it is completely isometric to
\(\mathbb C\).
\end{corollary}

\begin{proof}
This is Theorem~\ref{thm:integral-structure} with \(\lambda=1\).
\end{proof}

\begin{corollary}\label{cor:integer-jumps}
If \(V\neq 0\) is \(\lambda\)-nuclear in
\((\CI(\cdot,\cdot),\iota)\), then \(V\) is finite-dimensional and
\[
        \dim V=\iota(\id_V)\leq \lambda.
\]
In particular the admissible constants for non-zero spaces occur in
integer jumps.
\end{corollary}

\section{Consequences for the original statements}

The corrected form of \cite[Lemma 3.13]{DTZ} is
\[
        \iota(\id_{M_{n,k}})=nk.
\]
The corrected form of \cite[Theorem 3.14]{DTZ} is
\[
        M_{n,k}\hbox{ is }\lambda\hbox{-nuclear in }
        (\CI(\cdot,\cdot),\iota)
        \quad\Longleftrightarrow\quad
        nk\leq \lambda.
\]
Theorem~\ref{thm:integral-structure} gives the corresponding global
classification.  Thus the completely integral system has the same finite
rectangular building blocks as the completely nuclear system, but the
constant is the total dimension of those blocks.

\begin{thebibliography}{99}

\bibitem{BP}
D.~Blecher and V.~Paulsen,
\emph{Tensor products of operator spaces},
J. Funct. Anal. \textbf{99} (1991), 262--292.

\bibitem{B}
D.~Blecher,
\emph{Tensor products of operator spaces, II},
Canad. J. Math. \textbf{44} (1992), 75--90.

\bibitem{DTZ}
Z.~Dong, J.-C.~Tao and Y.-F.~Zhao,
\emph{Nuclearity and finite-representability in the system of completely
integral mapping spaces},
Acta Math. Sin. (Engl. Ser.) \textbf{40} (2024), 1197--1214.

\bibitem{ER-mapping}
E.~G. Effros and Z.-J. Ruan,
\emph{Mapping spaces and liftings for operator spaces},
Proc. London Math. Soc. (3) \textbf{69} (1994), 171--197.

\bibitem{EOR}
E.~G. Effros, N.~Ozawa and Z.-J. Ruan,
\emph{On injectivity and nuclearity for operator spaces},
Duke Math. J. \textbf{110} (2001), 489--521.

\bibitem{ER-book}
E.~G. Effros and Z.-J. Ruan,
\emph{Operator Spaces},
London Mathematical Society Monographs, New Series, 23,
Oxford University Press, New York, 2000.

\bibitem{P}
G.~Pisier,
\emph{Introduction to Operator Space Theory},
London Mathematical Society Lecture Note Series, 294,
Cambridge University Press, Cambridge, 2003.

\bibitem{R}
Z.-J.~Ruan,
\emph{Injectivity of operator spaces},
Trans. Amer. Math. Soc. \textbf{315} (1989), 89--104.

\end{thebibliography}

\end{document}
